\chapter{Introduction} \label{chap:intro}

\section*{}
% O primeiro capítulo da dissertação deve servir para apresentar o
% enquadramento e a moti\-va\-ção do trabalho e para identificar e
% definir os problemas que a dissertação aborda.
% Deve resumir as metodologias utilizadas no trabalho e termina
% apresentando um breve resumo de cada um dos capítulos
% posteriores.

% apresentar o Enquadramento:
%       MIR Visualization and Score-Alignemnt
% motivação do trabalho:

% identificar e definir os problemas que a dissertação aborda:


\section{Contextualization} \label{sec:context}
% What is MIR. Visualization in MIR. 
MIR is heavily reliant on data, and thus, visualization plays a crucial role; how and what we visualize in data can shape not only the conclusions we make but can also provide guidance on how to further progress. The problem of visualization in MIR is especially critical because music is very much subjective and abstract by virtue of being an artistic form of expression. MIR data visualization can span from audio representations, plots, to the various forms of musical scores. Visualization also plays an important role for annotation tasks providing annotators with relevant information to better aid their tasks. 
% N gosto muito disto, tenho de escrever e desenvolver isto melhor.

\section{Developed Project} \label{sec:proj}
%esta secção apresenta resumidamente o projeto.
For this thesis we developed MIRVisualScore, a tool that combines multiple state‑of‑the‑art and well established MIR resources into one single Python‑based object‑oriented (OO) tool. MIRVisualScore takes advantage of ScoreWarp, librosa, matplotlib, and BeatThis, in order to produce score aligned audio visualizations with overlaid algorithm outputs. 

\section{Motivation} \label{sec:goals}

% Apresenta a motivação e enumera os objetivos do trabalho terminando
% com um resumo das metodologias para a prossecução dos objetivos.

\section{Thesis Structure} \label{sec:struct}

